\chapter{Game Theory and Machine Learning Applications}
\label{sec:chapter11}

This chapter reviews theoretical and algorithmic optimization models for wireless MAC protocols, specifically game theory and machine learning models.

\section{Application of Game theory to scheduling}

A classical problem that has captured the attention of scientists and engineers from the networking community for decades is to design efficient techniques for fair allocation of scarce resources (like, spectrum, power, computing resource, and so on) among different contenders. 
%\cite{ResourceAllocation}. 
%The contenders can either be homogeneous or heterogeneous in nature. Similarly, both centralized as well as distributed algorithms can be applied in solving these problems.
%Such problems attempting efficient distribution of resources are interesting both from theoretical as well as application viewpoints. 
In particular, the use of game theory in addressing various resource allocation issues in wireless network (like congestion control, power control, medium access control, routing, edge computing, and so on) has increased many-folds over the past few years. 
However, recently the non-conventional game-theoretic models (like Minority Game) have attracted the attention of the networking community to overcome the limitations of conventional game models, that includes, 
%the following limitations of conventional game models resulted in increased focus on non-conventional game 
%he rationale behind choosing MG-based framework in our design is to address the following limitations of conventional game models while modelling large-scale networks: 
%The following limitations of applying conventional game models in designing game-based solutions for large-scale networks led us to choose MG as the underlying framework for our design:
%While designing game-based solutions for large-scale networks, several limitations in applying conventional game models have been identified. Limitations include \cite{Semasinghe16}:
%\begin{enumerate}
%[label=(\roman*)]
    (a) excessive signaling overhead related to data exchange between players,
    (b) high energy consumption due to computational complexity, 
    (c) assumption of complete rationality of players,
    (d) theoretical complexity in characterizing equilibrium set, and so on.
%\cite{SurveyGame2}
%\cite{GT_Book35}. 
%\section{Aim of our Research}
%\noindent
%In this thesis, we have concentrated upon three key networking problems for Next-Generation wireless network - (i) content distribution in Device-to-Device (D2D) communication, (ii) mobile Data Offloading via IEEE 802.11 WiFi network, and (iii) lightweight energy-efficient Medium Access Control (MAC) protocol design.
%for Internet-of-Things (IoT) applications. 
%In this research work, we have focused on design and analysis of distributed algorithms for the above mentioned resource allocation problems using Minority Game (MG) theoretic framework. 

\section{Overview of Contributions}
Contributions from researchers under the Mobile Computing Lab include:
\begin{itemize}
    \item \textbf{Bhaswar:} Game-theoretic formulations for channel selection, power control, and contention window adaptation.
    \item \textbf{Rohan, Lokesh, Indu:} Machine learning models (e.g., reinforcement learning) for dynamic spectrum access and adaptive rate selection.
\end{itemize}

\section{Game-Theoretic and ML-Based Resource Allocation}
We analyze how decentralized nodes interact strategically or adaptively learn to share scarce radio resources, summarizing Nash equilibria and RL-based MAC configurations.
